\chapter{Resumen}


\textbf{Resumen:} Este trabajo propone el desarrollo de un sistema inteligente para la recomendación vocacional de carreras 
del Instituto Politécnico Nacional en la Ciudad de México. La solución busca atender las limitaciones de herramientas 
tradicionales, que suelen ofrecer orientaciones estáticas y poco contextualizadas. El sistema integrará una plataforma web, 
una base de datos y una arquitectura modular basada en microservicios. Mediante cuestionarios de información personal, selección 
de área, habilidades y respuestas abiertas con metodología STAR, se recopilarán datos del usuario. Las técnicas de procesamiento 
de lenguaje natural permitirán extraer aptitudes de las respuestas, mientras que un modelo de clasificación supervisada 
relacionará dichas habilidades con programas académicos, escuelas y carreras del IPN. La implementación contempla servicios 
independientes, recuperación aumentada, comunicación orientada a eventos, PostgreSQL, Redis, despliegue en la nube y prácticas 
de integración continua. Se definirán pruebas unitarias, de integración y de evaluación para verificar el funcionamiento y los 
alcances de la propuesta.


\textbf{Palabras Clave:} orientación vocacional, aprendizaje automático, procesamiento de lenguaje natural, microservicios, clasificación supervisada, Instituto Politécnico Nacional.



%%%%%%fin del archivo

\endinput 